% 行内代码保持 C07 的 \texttt / \verb 语法，只统一字号与字体族。
\let\PortraitOriginalTexttt\texttt
\RenewDocumentCommand{\texttt}{m}{{\fontsize{8.1bp}{14bp}\selectfont\PortraitOriginalTexttt{#1}}}
\makeatletter
\def\verbatim@font{\normalfont\ttfamily\fontsize{8.1bp}{14bp}\selectfont}
\makeatother

% xcode 为试排候选，可按语言覆盖；源代码与缩进不变。
\newcommand\CodeFontSize{\fontsize{8.1bp}{12.5bp}\selectfont}
\definecolor{CodeNumber}{HTML}{999999}
\setminted{style=xcode,fontsize=\CodeFontSize,baselinestretch=1,
  linenos,firstnumber=1,numbersep=5bp,xleftmargin=22bp,
  tabsize=4,autogobble=false,mathescape=false,texcomments=false,
  showspaces=false,showtabs=false,breaklines=true,breakanywhere=true,
  breaksymbolleft={\textcolor{CodeNumber}{\ensuremath{\hookrightarrow}}},
  breaksymbolright={},breakindent=0pt}
\renewcommand\theFancyVerbLine{\textcolor{CodeNumber}{\ttfamily\fontsize{7bp}{9bp}\selectfont\arabic{FancyVerbLine}}}

% 与 HTML <mark> 相同的纯黄底色；零内边距，不改变短语尺寸。
% \colorbox 把短语作为不可拆分的行内盒子，局部分组防止间距设置外泄。
\newcommand{\Highlight}[1]{%
  \begingroup
  \setlength{\fboxsep}{0pt}%
  \colorbox[HTML]{FFFF00}{#1}%
  \endgroup
}

% 代码标题栏固定外观；代码字体、配色、行号与折行由 \setminted[语言] 控制。
\definecolor{InputCodeTitle}{HTML}{F2F2F2}
\ExplSyntaxOn
\str_new:N \l__inputcode_area_str
\str_new:N \l__inputcode_base_str
\str_new:N \l__inputcode_ext_str
\str_new:N \l__inputcode_title_str

\cs_new_protected:Npn \__inputcode_filename:n #1
  {
    \file_parse_full_name:nNNN {#1}
      \l__inputcode_area_str \l__inputcode_base_str \l__inputcode_ext_str
    \str_set:Nx \l__inputcode_title_str
      {\str_use:N \l__inputcode_base_str \str_use:N \l__inputcode_ext_str}
  }

\cs_new_protected:Npn \__inputcode_print_title:n #1
  {
    \str_map_inline:nn {#1}
      {
        \str_if_eq:nnTF {##1} { ~ }
          {\hspace{\fontdimen2\font}}
          {##1}
      }
  }

\cs_new_protected:Npn \__inputcode_box:nnn #1#2#3
  {
    \par
    \begin{tcolorbox}[
      enhanced,
      breakable,
      title={{\xeCJKsetup{CJKecglue={}}\__inputcode_print_title:n {#1}}},
      colback=white,
      colbacktitle=InputCodeTitle,
      colframe=black,
      coltitle=black,
      boxrule=0.4pt,
      titlerule=0.4pt,
      arc=0pt,
      left=1mm,right=1mm,top=1mm,bottom=1mm,
      lefttitle=1mm,righttitle=1mm,toptitle=0.5mm,bottomtitle=0.5mm,
      fonttitle={\normalfont\ttfamily\mdseries\upshape\fontsize{9bp}{14bp}\selectfont}
    ]
    \xeCJKsetup{CJKecglue={}}
    \inputminted[mathescape=false,texcomments=false,autogobble=false]{#2}{#3}
    \end{tcolorbox}
  }

\NewDocumentCommand{\InputCode}{o m m}
  {
    \IfNoValueTF{#1}
      {
        \__inputcode_filename:n {#3}
        \exp_args:NV \__inputcode_box:nnn \l__inputcode_title_str {#2} {#3}
      }
      {\__inputcode_box:nnn {#1} {#2} {#3}}
  }
\ExplSyntaxOff

% 可跨页的引用块：浅灰背景、左侧深灰竖线，无其余边框。
\definecolor{QuoteBlockBackground}{HTML}{F7F7F7}
\definecolor{QuoteBlockRule}{HTML}{777777}
\newtcolorbox{quoteblock}{
  enhanced,
  breakable,
  colback=QuoteBlockBackground,
  boxrule=0pt,
  borderline west={2pt}{0pt}{QuoteBlockRule},
  arc=0pt,
  boxsep=0pt,
  left=2mm,right=2mm,top=2mm,bottom=2mm
}

% Markdown 水平分隔线：按当前容器的 \linewidth 排版，不改变分页。
\definecolor{SeparatorRule}{HTML}{A0A0A0}
\NewDocumentCommand{\Separator}{}{%
  \par\addvspace{\baselineskip}%
  \noindent{\color{SeparatorRule}\rule{\linewidth}{0.4pt}}\par
  \addvspace{\baselineskip}%
}
